\chapter[Fundamentos de la Psicología infantil en entornos digitales]{Fundamentos de la Psicología infantil en entornos digitales}
\label{cp:Fundamentos-Psicologia}

\insertminitoc
\parindent1.5em 

\section{Desarrollo cognitivo y emocional en la infancia}

El desarrollo neurológico durante la infancia y la adolescencia se define por una plasticidad sináptica excepcional, lo que convierte a las funciones ejecutivas en procesos altamente susceptibles a los estímulos del entorno digital. Según \cite{clemente-suarez_digital_2024}, esta etapa se caracteriza por un refinamiento de las conexiones neuronales que puede verse alterado por la sobreestimulación de los dispositivos móviles. La investigación detalla que la exposición constante a interfaces de respuesta rápida puede generar una "atención fragmentada", dificultando que el menor desarrolle la capacidad de concentración profunda necesaria para tareas académicas y de resolución de problemas complejos. Esta fragmentación no solo afecta el rendimiento escolar, sino que altera la capacidad de introspección y autorregulación emocional a largo plazo \cite{limone_psychological_2021}. \\

Mientras que el uso moderado de tecnologías educativas puede ofrecer beneficios pedagógicos, la dependencia digital severa se ha correlacionado negativamente con la maduración de la corteza prefrontal y la integridad de la materia blanca \cite{joseph_impact_2022, clemente-suarez_digital_2024}. Esta región cerebral es la encargada del control de impulsos; un estudio que evaluó el comportamiento de menores tras periodos prolongados de uso digital encontró que aquellos con mayor tiempo de pantalla presentaban una menor densidad de materia blanca en áreas asociadas con el lenguaje y la alfabetización. Esto justifica la implementación de sistemas de supervisión que actúen como un andamiaje externo, apoyando al menor mientras su propio sistema de control inhibitorio termina de madurar biológicamente, evitando que la impulsividad digital se traduzca en conductas de riesgo permanentes \cite{borges_off_2025}. \\

Adicionalmente, la literatura científica ha explorado cómo el uso de dispositivos afecta los ritmos circadianos y, por ende, la salud mental. Mediante el análisis de hábitos de sueño, se ha determinado que la supresión de melatonina provocada por la luz azul no solo causa insomnio, sino que incrementa la vulnerabilidad a trastornos del estado de ánimo y reduce la resiliencia psicológica \cite{swider-cios_young_2023, page_psychological_2025}. Un menor cansado tiene menos recursos psicológicos para identificar y evitar riesgos en línea, como el ciberacoso o la manipulación. Por tanto, la supervisión técnica que propones responde a una necesidad fisiológica de proteger los ciclos de descanso, garantizando que el cerebro infantil cuente con los periodos de recuperación necesarios para un desarrollo emocional estable \cite{serra_smartphone_2021}. \\

Para visualizar el impacto multifactorial del entorno digital en el desarrollo biológico y psicológico, se presenta la siguiente tabla que sintetiza los hallazgos de autores clave:\\

\begin{table}[h!]
\centering
\caption{Efectos del entorno digital en el desarrollo infantil}
\begin{tabular}{|p{3.5cm}|p{4cm}|p{3.5cm}|p{3cm}|}
\hline
\textbf{Área de Impacto} & \textbf{Referencia Principal} & \textbf{Efecto Observado} & \textbf{Consecuencia Psicológica} \\ \hline
Estructura Cerebral & \cite{joseph_impact_2022, clemente-suarez_digital_2024} & Menor densidad de materia blanca en áreas prefrontales. & Déficit en control de impulsos y lenguaje. \\ \hline
Atención y Cognición & \cite{clemente-suarez_digital_2024, limone_psychological_2021} & Fragmentación de la atención por respuesta rápida. & Dificultad para el aprendizaje profundo. \\ \hline
Salud Circadiana & \cite{swider-cios_young_2023, serra_smartphone_2021} & Supresión de melatonina por luz azul. & Insomnio y labilidad emocional. \\ \hline
Bienestar Emocional & \cite{page_psychological_2025, borges_off_2025} & Trauma acumulativo por victimización. & Vulnerabilidad a largo plazo. \\ \hline
\end{tabular}
\end{table}

\section{Impacto de la frecuencia de uso en la primera infancia}

La sobreexposición a dispositivos móviles durante los primeros años de vida genera alteraciones conductuales específicas que difieren de las observadas en adolescentes. Según un estudio de revisión sistemática de \cite{pluas_frequent_2024}, los altos niveles de uso de dispositivos móviles a la temprana edad de 5 años están fuertemente asociados con hiperactividad, dificultades graves de atención y problemas de concentración. El estudio alerta que incluso en niños de 18 meses, el uso frecuente obstaculiza el desarrollo de las redes de atención, ya que los infantes son altamente susceptibles a las experiencias mediáticas externas que interrumpen los ritmos de aprendizaje natural y el juego libre no estructurado.\\

Estas alteraciones cognitivas se ven exacerbadas cuando el uso ocurre sin supervisión adulta. Investigaciones transaccionales en ciudades con alta penetración digital, como el estudio realizado por \cite{kakon_supervised_2025} en Bangladesh, revelan que la falta de supervisión del tiempo de pantalla se correlaciona con un aumento en quejas somáticas como dolores de cabeza, problemas de visión y dolores abdominales. En su investigación con escolares, observaron que los niños bajo supervisión parental mantenían promedios de sueño más saludables, lo que sugiere que la ausencia de límites digitales se traduce casi inmediatamente en deficiencias fisiológicas que impactan el rendimiento cognitivo al día siguiente.\\

El uso temprano de tecnología también altera la forma en que los niños aprenden a gestionar sus propias emociones. Cuando los padres o tutores utilizan dispositivos digitales como un "chupete electrónico" para calmar rabietas o episodios de llanto, se interrumpe el proceso de aprendizaje de la autorregulación emocional \cite{taylor_does_2024}. La literatura psicológica sostiene que, si bien el co-uso (adulto y niño interactuando juntos con el dispositivo) puede tener una asociación positiva menor en el aprendizaje temprano, el uso solitario y no guiado priva al menor de las interacciones cara a cara esenciales para desarrollar la empatía y la resiliencia ante la frustración. \\

Por tanto, el diseño de herramientas tecnológicas de protección infantil debe contemplar métricas que vayan más allá de la restricción horaria. Es imperativo medir indicadores de uso solitario, uso antes de dormir o uso durante rutinas familiares (como las comidas) para identificar patrones de "uso problemático", ya que estos comportamientos de riesgo son más predecibles de futuros trastornos de salud mental que la simple exposición cuantitativa a las pantallas \cite{selak_problematic_2025}.\\

\section{Influencia del entorno digital en la conducta infantil}

El entorno digital funciona como un ecosistema diseñado bajo principios de psicología conductual, utilizando mecanismos de recompensa intermitente que influyen directamente en la conducta del niño. Se ha documentado que el Uso Problemático de Internet (UPI) correlaciona con problemas de internalización y externalización, afectando la forma en que los menores procesan sus emociones negativas \cite{bohnert_emerging_2021}. Un análisis longitudinal demostró que el uso excesivo de tecnología está vinculado con una disminución de la conducta prosocial y un aumento de la agresividad, especialmente cuando el contenido digital desplaza las interacciones cara a cara con pares y familiares, creando un vacío de socialización real \cite{masri-zada_impact_2025}. \\

El bienestar socioemocional del menor sufre alteraciones medibles cuando el tiempo de pantalla excede las 3 horas diarias. Investigaciones recientes sugieren que este umbral marca el inicio de una degradación en la satisfacción vital y un aumento en la ansiedad social, exacerbado por la comparación constante en redes sociales \cite{soriano-molina_association_2025, limone_psychological_2021, celestino_smart_2025}. En este punto, aparece el fenómeno del $FoMO$ (miedo a perderse algo), que actúa como un motor de ansiedad constante. Los menores sienten la necesidad de estar permanentemente conectados para validar su estatus social, lo que los empuja a participar en dinámicas de grupo potencialmente nocivas \cite{yusuf_lets_2023, holmarsdottir_integrative_2025}. \\

Asimismo, la "tecnoferencia" o la intrusión de la tecnología en las relaciones humanas erosiona el vínculo afectivo básico. Un estudio integrador sobre la vida digital infantil destaca que cuando los padres o los hijos están sumergidos en sus dispositivos durante momentos de convivencia, se debilita el aprendizaje de la empatía \cite{selak_problematic_2025, soriano-molina_association_2025}. Esta desconexión facilita que el menor busque refugio emocional en comunidades virtuales con desconocidos. Tu propuesta de una bitácora de eventos busca precisamente revertir este efecto, proporcionando un punto de encuentro técnico que obliga a retomar la conversación familiar \cite{stoilova_parental_2024, theopilus_parental_2025}. \\

\section{Vulnerabilidad psicológica, ciberacoso y riesgos de victimización}

La vulnerabilidad en el ciberespacio es una extensión de las vulnerabilidades preexistentes en el entorno físico del menor. Investigaciones demostraron que aquellos con baja autoestima o falta de supervisión presentan una predisposición significativamente mayor a ser víctimas de \textit{grooming} o estafas \cite{el-asam_psychological_2022, ayyash_cybersecurity_2024}. El distrés psicológico actúa como un catalizador: un menor que se siente solo es más propenso a compartir información privada a cambio de atención externa, una dinámica que los agresores explotan mediante tácticas de manipulación sofisticadas \cite{razi_sliding_2023, page_psychological_2025}. \\

Esta vulnerabilidad se agrava por el fenómeno de la desinhibición digital, donde el menor siente que las consecuencias en línea no son reales. Una revisión de estudios globales resalta que la falta de presencia física de un adulto facilita que el niño sea tanto víctima como agresor accidental \cite{zhu_cyberbullying_2021, yusuf_lets_2023}. La investigación identifica que la mediación técnica es vital, ya que el 87\% de los padres considera la seguridad digital una prioridad, pero solo la mitad aplica controles efectivos \cite{ayyash_cybersecurity_2024, aznar-martinez_uso_2024}. Aquí radica la importancia de la bitácora de evidencias: permite detectar el problema sin esperar a que el menor venza la barrera del silencio. \\

El ciberacoso (\textit{cyberbullying}) se destaca como uno de los riesgos de victimización más prevalentes, afectando a la juventud en un nivel global, sin distinción de fronteras \cite{craig_social_2020}. La "cara oscura de la era digital" ha facilitado que el acoso sea omnipresente, siguiendo al menor hasta su propio dormitorio a través del teléfono móvil, lo que genera cuadros de ansiedad severa, depresión y, en casos extremos, ideación suicida \cite{wiederhold_dark_2024}. Frente a esta amenaza, los sistemas de protección deben evolucionar desde el simple bloqueo hasta el análisis profundo de texto, para captar las agresiones antes de que se consoliden en traumas irreversibles.\\

\begin{table}[h!]
\centering
\caption{Factores de vulnerabilidad y riesgos asociados en menores}
\begin{tabular}{|p{3.5cm}|p{4cm}|p{4cm}|p{3cm}|}
\hline
\textbf{Factor de Riesgo} & \textbf{Referencia Principal} & \textbf{Descripción del Fenómeno} & \textbf{Riesgo Resultante} \\ \hline
Baja Autoestima / Soledad & \cite{el-asam_psychological_2022, page_psychological_2025} & Búsqueda de validación en extraños. & Grooming y extorsión sexual. \\ \hline
Brecha de Supervisión & \cite{ayyash_cybersecurity_2024, rojas_technical_2025} & Padres conscientes pero inactivos técnicamente. & Exposición a contenido sexual y fraudes. \\ \hline
Desinhibición Digital & \cite{zhu_cyberbullying_2021, yusuf_lets_2023} & Sensación de anonimato e impunidad. & Ciberacoso y difusión de datos. \\ \hline
Uso Excesivo (>3h) & \cite{soriano-molina_association_2025, bohnert_emerging_2021} & Desplazamiento de actividades saludables. & UPI y ansiedad social. \\ \hline
\end{tabular}
\end{table}

\section{Resiliencia digital y el enfoque de Desarrollo Positivo Juvenil (PYD)}

El paradigma actual de la psicología digital se aleja de los enfoques basados en el miedo para enfocarse en la "resiliencia digital" y el Desarrollo Positivo Juvenil (PYD, por sus siglas en inglés). Revisiones sistemáticas han demostrado que potenciar los factores de protección internos del adolescente, como la autorregulación emocional y el pensamiento crítico, funciona como un escudo contra las conductas de riesgo en línea como la exposición a pornografía, el \textit{sexting} o la adicción a los juegos \cite{martin-barrado_association_2025}. Este enfoque busca que el menor desarrolle las competencias necesarias para navegar de forma segura por sí mismo, transformándolo en un sujeto empoderado más que en una víctima potencial.\\

Las investigaciones determinaron que la mediación habilitadora es la que mejores resultados ofrece a largo plazo para fomentar esta resiliencia \cite{livingstone_maximizing_2017, stoilova_parental_2024}. Este estilo de mediación requiere que el padre esté presente en la vida digital del hijo, guiándolo y proporcionándole herramientas para identificar engaños, lo cual es solo posible si el tutor tiene acceso a información veraz sobre la actividad del menor \cite{zgambo_effect_2025, allison_parental_2025}. Estudios indican que las prácticas parentales que fomentan el co-uso o la supervisión activa reducen drásticamente la incidencia de uso problemático en adolescentes tempranos \cite{nagata_associations_2025}.\\

La aceptación de las herramientas de supervisión por parte de los niños depende directamente de la transparencia y de cómo la aplicación fomenta su propia agencia. Estudios de co-diseño revelaron que los menores prefieren sistemas que promuevan la comunicación sobre aquellos que bloquean aplicaciones sin dar justificación alguna \cite{mcnally_co-designing_2018, holmarsdottir_integrative_2025}. Los niños valoran la utilidad de la supervisión cuando se les permite revisar la evidencia junto a sus padres, entendiendo el proceso como una guía en lugar de un control autoritario. Este proceso de co-regulación fortalece el vínculo afectivo y reduce la reactancia psicológica, transformando la aplicación de monitoreo en un facilitador educativo \cite{chowdhury_landscape_2025, liverpool_engaging_2020}. \\

Al automatizar la detección de incidentes mediante \gls{ia}, el proyecto reduce la carga cognitiva y la fatiga parental, asegurando que el tiempo de interacción familiar se dedique al apoyo emocional constructivo y a la discusión ética \cite{stoilova_parental_2024, ayyash_cybersecurity_2024, koulaxidis_improving_2022}. La meta final es la autonomía segura. Al usar una bitácora de evidencias protegida y local, se crea un historial de aprendizaje donde el menor puede ver sus propios errores bajo una luz constructiva, equilibrando la innegable necesidad de protección infantil con el derecho a la libertad necesaria para el crecimiento sociocognitivo en el siglo XXI \cite{agbana_ethical_2025, balapour_mobile_2020}. \\

\section{Alineación con los Objetivos de Desarrollo Sostenible (ODS)}

El desarrollo y la fundamentación de este sistema de supervisión inteligente se enmarcan directamente en la Agenda 2030, contribuyendo al cumplimiento de los Objetivos de Desarrollo Sostenible (ODS) promovidos por la Organización de las Naciones Unidas (ONU). De manera específica, el modelo tecnológico y psicológico propuesto en esta investigación incide en los siguientes objetivos:

\textbf{ODS 3: Salud y bienestar.} El uso problemático de internet, la disrupción de los ritmos circadianos y la victimización por ciberacoso tienen un impacto directo y negativo en la salud mental infantil, generando cuadros de ansiedad, aislamiento social y depresión. Al proporcionar una herramienta preventiva que detecta patrones de riesgo antes de que se consoliden en traumas, esta tesis contribuye directamente a garantizar una vida sana y promover el bienestar psicológico y emocional de los menores, protegiendo su desarrollo cognitivo en una etapa crítica.

\textbf{ODS 4: Educación de calidad.} La alfabetización y la resiliencia digital son componentes ineludibles de la educación integral en el siglo XXI. Al integrar un enfoque de Desarrollo Positivo Juvenil y transformar la bitácora de evidencias en un instrumento pedagógico para la familia, el proyecto fomenta un ambiente de aprendizaje seguro y no violento. El sistema dota a los padres de herramientas para educar a sus hijos en el uso ético y crítico de la tecnología, asegurando que los entornos de aprendizaje digital sean inclusivos y efectivos.

\textbf{ODS 16: Paz, justicia e instituciones sólidas.} Este objetivo establece en su Meta 16.2 la necesidad imperativa de poner fin al maltrato, la explotación, la trata y todas las formas de violencia contra los niños. Al implementar algoritmos de \gls{ia} capaces de interceptar intentos de \textit{grooming}, extorsión sexual infantil (OCSEA) y violencia verbal en plataformas móviles, la herramienta actúa como un mecanismo directo de justicia y protección. Además, al procesar los datos localmente y garantizar la privacidad, la solución fortalece los derechos humanos del menor, promoviendo instituciones familiares más informadas y sociedades más justas en el ecosistema cibernético.